\section{Four Pell roots: fourteen integers at degrees 540 and 568}
\label{sec:crt14}

This section gives two ordinary polynomials whose fourteen coordinates all
range over the unrestricted integers. The first has exact degree $540$ and
$160$ live binary arithmetic operations; the second has degree $568$ and
$147$ operations. The fixed coefficient tuple has $615\,883$ raw minimal
unsigned bytes, with the original six integer payloads retained. Operations
count shared binary addition, subtraction, and multiplication, including
multiplication by a fixed constant. A register is a polynomial abbreviation,
not an additional quantified coordinate.

\subsection{Literal definition}

The variables, in circuit order, are
\[
 j,C,z,h,r,v,w,T,c,n,t,p_2,p_3,p_1\in\mathbb Z.
\]
Use the six fixed integers of the retained source-34/source-71 tuple:
\[
 \beta=2^{1064921}-1,\quad K_0,M_C,M_F,\quad
 A_{\mathrm c}=2^{1113},\quad B_{\mathrm c}=2^{666104}.
\]
Here $M_F$ is already shifted by $\beta$. The exact integer bytes, finite
machine-and-tile recipe, original licenses, and ordinary calibration texts
are included in the component's \texttt{upstream/} directory. This tuple is
larger than some other native tuples in ProveIt; the present determinant
scales are proved for this whole tuple and are not mixed with another one.

Define
\begin{align*}
 J&=6j+1,&q&=\beta J+1,&W&=12w+4,\\
 X&=qW,&Y&=qv,&E&=XY,\\
 a&=E+Y,&H&=4a+3,&\Delta&=a^2+H.
\end{align*}
Then define the transport and rank registers
\begin{align*}
 F&=(K_0+W)C-z(q-1),\\
 R&=(q(q-F)-C)(q^2-1)+(M_C+qM_F)J.
\end{align*}
Continue with
\begin{align*}
 k&=R+1+hE,&\eta&=c-Yk,&\zeta&=k-\eta,\\
 i&=2r+1,&I&=ic^2,&D_f&=1+\Delta I^2,\\
 Q&=\Delta^2I^2,&\delta&=Q-1,&V&=D_f(cT-R)-c,\\
 C_y&=QV^2-1,&\ell&=XY^2,&L&=\ell k,\\
 D_\tau&=1+\ell(\ell+1)k^2,&D_m&=1+\Delta c^2,\\
 b_0&=X+ac,&F_0&=(a+2)I.
\end{align*}
Put
\begin{align*}
 \operatorname{Pos}(u_1,u_2,u_3)&=1+u_1^2+u_2^2+u_3(u_3+1),\\
 G&=(A_{\mathrm c}F)(B_{\mathrm c}C)
       -(q-A_{\mathrm c}F-B_{\mathrm c}C)^2,\\
 B&=v^2G\eta\zeta-(t^2-v^2G-\eta\zeta+1)^2
       -\operatorname{Pos}(p_1,p_2,p_3),\qquad S=B^2.
\end{align*}
Let
\[
 \mathcal F_{16}(x;A_1,A_2,A_3,A_4)
 =\prod_{\varepsilon\in\{-1,1\}^4}
       \left(x-\sum_{j=1}^4\varepsilon_j\sqrt{A_j}\right).
\]
Sign invariance makes this an integer polynomial in its displayed
arguments. Its exact $495$ nonzero coefficients are verified against the
independent product of the sixteen factors.

For the centered equation define
\begin{align*}
 M&=4q^2,&U&=MH,\\
 Z&=U(n+2V+2L+k+2F_0)+2Mb_0+2H(v-q^2),\\
 N&=\delta^8\mathcal F_{16}
   \left(Z;4U^2D_f,4U^2D_\tau,4M^2D_m,\frac{4U^2C_y}{\delta}\right).
\end{align*}
\[\boxed{P_{540}=N+S+NS=0.}\]
The denominator is cleared symbolically; the literal evaluator contains
only integer polynomial operations. For the shorter companion use instead
\begin{align*}
 M&=2q^2,&U&=MH,\\
 Z&=U(n+V)+Mb_0+H(v-q^2),\\
 N&=\delta^8\mathcal F_{16}
 \left(Z;U^2D_f,U^2D_\tau,M^2D_m,\frac{U^2C_y}{\delta}\right).
\end{align*}
\[\boxed{P_{568}=N+S+NS=0.}\]
All other registers remain unchanged.

\subsection{The two guard facts}

For integers $x,y,t$ define
\[
 \operatorname{Pair}(x,y,t)=xy-(t^2-x-y+1)^2.
\]
Then
\[
 (\exists t\in\mathbb Z)\;\operatorname{Pair}(x,y,t)>0
 \quad\Longleftrightarrow\quad x>0\ \text{and}\ y>0.
\]
The forward direction follows because positive Pair forces $xy>0$, whereas
for $x=-u,y=-v$ with $u,v>0$ its subtracted square is at least
$(u+v+1)^2>uv$. For the reverse direction set $K=x+y-1$ and choose a
nearest integer square $t^2$. Between consecutive squares $m^2,(m+1)^2$
the error is at most $m$, so $(t^2-K)^2<K\le xy$; square endpoints,
including $K=1$, have zero error.

The polynomial Pos represents precisely the positive integers. Positivity
is immediate; surjectivity follows from the three-square theorem applied
to $4n-3$, with one odd and two even square coordinates. Consequently,
$B=0$ forces $v\ne0$, $G>0$, and $\eta\zeta>0$. Conversely these three
conditions admit private witnesses making $B=0$.

Since $S=B^2\ge0$, the equation
\[
 P=(1+S)(1+N)-1=0
\]
forces $S=N=0$: $1+S$ is a positive integer divisor of $1$. No sign
hypothesis about the norm $N$ is used.

\subsection{Rational extraction before the cubic scale is known}

The front end gives $q\equiv2\pmod6$ and $|q|\ge16$. Also
$W=4(3w+1)\ne0$, hence $x=|X|\ge64$. When $v=0$ we have $a=0$,
$H=3$, $\Delta=3$; when $v\ne0$, $m=|Y|\ge16$ and $|a|\ge1008$.
Thus $H\ne0$ and $\Delta\ge3$ on the entire integer domain.
The three integer radicands $D_m,D_f,D_\tau$ are positive, the last by
$\ell(\ell+1)\ge0$.

If $c=0$, then $Q=0$, $\delta=-1$, and $C_y/\delta=1$. If $c\ne0$,
then $Q\ge9$, $D_f>|c|+1$, and the integer $V=D_f(cT-R)-c$ has
$|V|\ge1$. Therefore
\[
 \frac{C_y}{\delta}=V^2+\frac{V^2-1}{Q-1}\ge1.
\]
In particular $\delta$ never vanishes, including all exceptional branches.

At a polynomial zero, the guard supplies $v\ne0$, $G>0$, and
$\eta\zeta>0$. Write $K=|k|$, $C_0=|c|$. Then $K\ge2$ and
$0<\eta/k<1$, so
\[
 K(m-1)<C_0<K(m+1),\qquad C_0>30.
\]
Let $d,f,\tau,y$ be the positive square roots of
$D_m,D_f,D_\tau,C_y/\delta$, respectively. They satisfy
\[
 \frac{d}{|H|}<\frac{C_0}{2},\qquad
 \tau>2C_0,\qquad f>2\tau.
\]
For the first inequality use $H^2-4\Delta=12a^2+8a-3$.
For the second use
$K(xm^2-1)<\tau<K(xm^2+1)$.
For the third,
\[
 d>((x-1)m-3)K(m-1)>\frac34\tau,
 \qquad f>\frac{C_0}{2}d>15d.
\]
The middle comparison follows by expanding
$m^2(x-4)-4m(x+2)+9>0$ for $x\ge64,m\ge16$.
Also $f>3C_0^2$ and $Q>f^2$.

If $cT-R=0$, then $C_0<y<C_0+1$. Otherwise
$y>|V|\ge f^2-C_0>2f$. Hence either
\[
 f,\tau,y,d/|H|\quad\text{or}\quad y,f,\tau,d/|H|
\]
is a strictly superincreasing descending list: every entry exceeds the
sum of all later entries. For example,
$y+d/|H|<1.5C_0+1<2C_0<\tau$ in the first case.
The deleted source amplifier $4H+1$ is unnecessary because the guard root
is now a separate explicit coordinate, leaving these four roots to separate.

A rational signed sum of positive roots with rational squares and strictly
superincreasing magnitudes makes every root rational. Indeed, a nontrivial
Galois automorphism changing at least one radical sign would yield a
nonempty vanishing signed subsum after subtraction, contradicting the
largest-term inequality. This argument allows dependent square classes.
At $N=0$ one conjugate factor vanishes; division by $2U$ in the centered
case, or $U$ in the companion, gives exactly such a rational sum.
Thus $d,f,\tau$ are integers.

The auxiliary rational $y$ is also integral. For any integers $Q>1,V$,
\[
 QV^2-(Q-1)y^2=1,\quad y\in\mathbb Q
 \quad\Longrightarrow\quad y\in\mathbb Z.
\]
If $y=a/b$ is reduced, then $b^2\mid Q-1$. The transformation
\[
 V'=(2Q-1)V-2(Q-1)y,\qquad
 y'=(2Q-1)y-2QV
\]
preserves the equation and denominator, keeps $V'$ integral, and for
$V>1$ gives $-V<V'<V$. Here $y>V$ and $(Q-1)y<QV$ give the two
strict inequalities. The equation excludes $V'=0$. Descent reaches
$|V|=1$, where $|y|=1$, proving the claim.

\subsection{The congruences and the consistency equivalence}

Absorb the conjugate signs into integer roots. In the centered case their
relation is
\[
 Z=2Md+2U(f+\tau+y).
\]
Since $H\equiv3\pmod q$, $q\equiv2\pmod3$, and $H$ is odd,
$\gcd(H,4q^2)=1$. Reduction modulo $M=4q^2$ yields
$2q^2\mid v-q^2$, and reduction modulo $H$ yields $H\mid d-b_0$.
Therefore
\[
 v=q^2(2v_{\rm old}+1),\quad
 Y=(2v_{\rm old}+1)q^3,\quad
 \gamma=(d-b_0)/H\in\mathbb Z.
\]
The companion with $M=2q^2$ gives the same conclusions without the
factors $2$. No cubic-scale promise was used in the preceding extraction.

These roots, the restored scale, and the two strict guards give the retained
parent's normalized signed kernel. Its original strong variables can be
restored explicitly:
\begin{align*}
 I_{\rm old}&=\Delta i c^2,& r_{\rm old}&=(\Delta i-1)/2,\\
 o_{\rm old}&=f(cT-R),&j_{\rm old}&=D_fT-1-R\Delta i^2c^3.
\end{align*}
The quotient is integral because $\Delta i$ is odd. Substitution gives
$I_{\rm old}^2=\Delta(f^2-1)=Q$ and
$fo_{\rm old}-c=cj_{\rm old}-R=V$. Private old guard witnesses are
regenerated from $G>0,\eta\zeta>0$.

Conversely, take a normalized source witness, set
$v=q^2(2v_{\rm old}+1)$, and choose the private guard witnesses by Pair
and Pos. For the centered equation choose
\[
 n=2(\gamma+f+\tau+y-V-L-F_0)-k
        -\frac{v-q^2}{2q^2}.
\]
For the companion choose
\[
 n=\gamma+f+\tau+y-V-\frac{v-q^2}{2q^2}.
\]
Both are integers and make the selected conjugate factor vanish. Every
quantity is chosen before $n$, so the construction is not circular.

We have proved an exact two-way existence reduction to the normalized
kernel. Composing with the parent's explicitly identified ordinary native
compiler and K215/refutation calibration yields, for $d=540,568$,
\[
 \mathrm{ZFC}\vdash
 \left(\neg\exists\boldsymbol x\in\mathbb Z^{14}
              \ P_d(\boldsymbol x)=0\right)
 \longleftrightarrow\operatorname{Con}(\mathrm{ZFC}).
\]
This imports ordinary mathematical theorems and their stated canonical
completeness. It does not promote finite arithmetic checks into an unbounded
semantic proof, or supply the separately open encoded-ZFC proof artifact.

\subsection{Exact degrees, including cancellation}

Treat the fixed coefficients as degree zero. For leading parts set
\begin{align*}
 q_0&=6\beta j,&W_0&=12w,&X_0&=q_0W_0,&Y_0&=q_0v,\\
 a_0&=X_0Y_0,&i_0&=2r,&R_0&=q_0^4z-q_0^3W_0C,\\
 k_0&=R_0+ha_0,&\eta_0&=-Y_0k_0,&Q_0&=a_0^4i_0^2c^4,\\
 V_0&=-a_0^2i_0^2c^4R_0,&L_0&=X_0Y_0^2k_0,&F_{00}&=a_0i_0c^2.
\end{align*}
The degrees of $\delta,U,V,L,F_0,\eta$ are $22,6,19,11,7,7$.
The guard has top part $-\eta_0^4$, so $S$ has degree $56$ and
leading part $\eta_0^8$.

For the centered equation use approximate roots
\[
 r_y=2UV,\quad r_\tau=U(2L+k),\quad
 r_f=2UF_0,\quad r_d=2Mb_0,
 \quad s=Un+2H(v-q^2).
\]
The center is their sum plus $s$. The four exact square defects are
\begin{align*}
 4U^2C_y/\delta-r_y^2&=4U^2(V^2-1)/\delta,\\
 4U^2D_\tau-r_\tau^2&=U^2(4-k^2),\\
 4U^2D_f-r_f^2&=4U^2(1-I^2),\\
 4M^2D_m-r_d^2&=4M^2(1+Hc^2-2acX-X^2).
\end{align*}
The approximate root degrees are $25,17,13,7$, and the corresponding
formal square-root correction degrees are at most $3,5,5,4$.
To justify this calculation, scale every variable by an indeterminate $u$
and work over the Laurent-series field of the rational function field in
the original coordinates. Each radicand is an approximate-root square
times $1$ plus negative-degree terms, so the binomial series supplies its
formal root with that leading sign. The delivered polynomial uses no series.

Eight conjugate factors have degree $25$, four degree $17$, two degree
$13$, and the last two degree $7$. For the final two, the leading terms
are $U_0n$ and $U_0(n+c)$, where $U_0=16q_0^2a_0$.
The product therefore has degree $308$, and the clearing factor
$\delta^8$ adds $176$. Thus $\deg N=484$, and
\[
 (P_{540})_{540}
 =2^{28}U_0^{16}Q_0^8V_0^8L_0^4F_{00}^2\eta_0^8 n(n+c).
\]
In particular the actual fixed-coefficient slice satisfies
\[
 [j^{240}w^{82}v^{94}h^{12}r^{34}c^{68}z^8n^2]P_{540}
   =2^{530}3^{322}\beta^{240}>0.
\]
The selected exponents sum to $540$, proving exact degree.

For the simpler companion only the auxiliary root is centered. Eight
factors have degree $25$ and eight have degree $17$; the norm has degree
$8\cdot25+8\cdot17+8\cdot22=512$. Now $U_0=8q_0^2a_0$, and
\[
 (P_{568})_{568}=2^8U_0^{16}Q_0^8V_0^8L_0^8\eta_0^8.
\]
Its selected coefficient is
\[
 [j^{256}w^{88}v^{104}h^{16}r^{32}c^{64}z^8]P_{568}
 =2^{520}3^{344}\beta^{256}>0.
\]
The exact audits check these entire leading forms, all four square-defect
identities, all $495$ coefficients of the generic norm, all six fixed
integer payloads against their finite reconstruction, and both literal
DAGs. Finite modular comparisons are additionally retained as transcription
diagnostics; the unbounded logical argument is the one above.
